\chapter{Spalart-Allmaras Solver Verification Details}
\label{app:B_sa_solver_verification}

This appendix supports Chapter~\ref{ch:sa_verification}. It contains concise FEniCS setup tables, the mesh-independence studies used for the channel benchmark, and the selected FEniCS and Ansys mesh information used for the U-bend verification case. For traceability, the active configuration file is named once near each FEniCS setup; the tables retain only the governing geometry, boundary-condition, and solver settings rather than duplicating the full Python configuration. The backward-facing-step material is placed in Appendix~\ref{app:A_backstep_benchmark}, because it supports the Chapter~\ref{ch:turbulence_modeling} turbulence-model discussion rather than the Chapter~\ref{ch:sa_verification} forward-solver verification.

Unless stated otherwise, all grid-independence relative differences in this appendix are reported as positive percentage magnitudes. For a mesh-dependent metric $\phi$, the coarse--medium difference is computed as $|\phi_{\mathrm{coarse}}-\phi_{\mathrm{medium}}|/|\phi_{\mathrm{coarse}}|\times100\%$, while the medium--fine difference is computed as $|\phi_{\mathrm{medium}}-\phi_{\mathrm{fine}}|/|\phi_{\mathrm{fine}}|\times100\%$.

\section{Verification Benchmark 1: Channel Flow}
\label{sec:app_A_channel}

\subsection{Physical and Geometric Parameters}
\label{subsec:app_A_channel_physical}
The first verification benchmark consists of a pressure-driven periodic 2D channel. To align with fundamental turbulence research conventions, the characteristic length is defined as the full channel height ($H$), and the Reynolds number ($Re_H$) is computed strictly based on this geometric parameter rather than the hydraulic diameter. The active FEniCS setup is summarized from \texttt{ConfigChannel\_SpalartAllmaras\_Steady.py}, and the key physical setup is detailed in Table~\ref{tab:app_A_channel_physical}.

\begin{table}[htbp]
	\centering
	\footnotesize
	\caption{Geometric, physical, and boundary-condition parameters for the periodic channel flow benchmark.}
	\label{tab:app_A_channel_physical}
	\renewcommand{\arraystretch}{1.2}
	\begin{tabularx}{\textwidth}{l X}
		\hline
		\textbf{Parameter} & \textbf{Channel flow setting} \\
		\hline
		Benchmark type & Pressure-driven, streamwise-periodic 2D channel \\
		Domain dimensions & Length $L=2.0$ m; full channel height $H=2.0$ m \\
		Characteristic length & Full channel height $H=2.0$ m \\
		Reference velocity & $U_{\mathrm{ref}}=18.5$ m/s, utilized to define the target Reynolds number \\
		Kinematic viscosity & $\nu = 1.81818 \times 10^{-3}$ m$^2$/s \\
		Reynolds number & $Re_H \approx 2.03 \times 10^4$ \\
			Body force & $\mathbf{f}=(1.0,0)$ m/s$^2$ in the kinematic-pressure convention \\
			Driving condition & Streamwise body-force mode, equivalent to a pressure drop of $2$ Pa over the periodic length $L=2$ m \\
			Velocity boundary conditions & Periodic velocity between inlet/outlet planes; no-slip on top and bottom walls \\
			Pressure boundary conditions & Periodic pressure field with no inlet/outlet pressure Dirichlet condition; pressure mean normalized during IPCS \\
			SA boundary conditions & Periodic $\tilde{\nu}$ between inlet/outlet planes; $\tilde{\nu}=0$ at walls \\
			Initial fields & $\mathbf{u}_0=(18.5,0)$ m/s, $p_0=0$ Pa, $\tilde{\nu}_0=5.0 \times 10^{-3}$ m$^2$/s \\
		Verification mesh & \texttt{Fine\_FEniCS}; first layer $\Delta y_1 = 1.709 \times 10^{-3}$ m \\
		\hline
	\end{tabularx}
\end{table}

\subsection{Numerical Simulation Parameters}
\label{subsec:app_A_channel_numerical}
The channel flow was resolved using an Incremental Pressure Correction Scheme (IPCS) embedded within a Picard fixed-point iteration loop. The key numerical solver settings, convergence tolerances, and relaxation factors used in FEniCS are summarized in Table~\ref{tab:app_A_channel_numerical}.

\begin{table}[htbp]
	\centering
	\scriptsize
	\caption{Numerical simulation parameters for the FEniCS steady Spalart-Allmaras channel benchmark.}
	\label{tab:app_A_channel_numerical}
	\renewcommand{\arraystretch}{1.2}
	\begin{tabularx}{\textwidth}{l c X}
		\hline
		\textbf{Parameter} & \textbf{Symbol} & \textbf{Setting} \\
		\hline
		Velocity space & $\mathbf{u}$ & Vector CG2, periodic in streamwise direction \\
		Pressure space & $p$ & Scalar CG1 \\
		SA working-variable space & $\tilde{\nu}$ & Scalar CG1, periodic in streamwise direction \\
		Quadrature degree & -- & $2$ \\
			Maximum outer Picard steps & $N_{\mathrm{Picard,max}}$ & $400$ \\
			SA solves per Picard step & -- & $1$ \\
			Outer velocity tolerance & $\epsilon_{\mathbf{u}}$ & $1.0 \times 10^{-5}$ \\
			Outer pressure tolerance & $\epsilon_p$ & $1.0 \times 10^{-6}$ \\
			Outer SA tolerance & $\epsilon_{\tilde{\nu}}$ & $1.0 \times 10^{-6}$ \\
			Outer diagnostic convergence & -- & Midline bulk velocity window of $100$ Picard steps with relative tolerance $1.0\times10^{-3}$ \\
			IPCS pseudo-time step & $\Delta t$ & $5.0 \times 10^{-3}$ s \\
			Maximum IPCS steps per Picard step & -- & $500$ \\
			IPCS velocity tolerance & -- & $1.0 \times 10^{-6}$ \\
			IPCS pressure tolerance & -- & $1.0 \times 10^{-6}$; pressure convergence is diagnostic in body-force mode \\
			IPCS pressure normalization & -- & Enabled; the pressure correction is mean-normalized \\
		Velocity relaxation factor & $\alpha_u$ & $0.3$ \\
		Pressure relaxation factor & $\alpha_p$ & $0.3$ \\
		SA relaxation factor & $\alpha_{\tilde{\nu}}$ & $0.3$ \\
		SA lower bound & $\tilde{\nu}_{\min}$ & $1.0 \times 10^{-12}$ \\
		Wall-distance equation & -- & Relaxed Eikonal, $\sigma_w=0.1$, $G_0=20$ \\
		Wall-distance Newton settings & -- & Relative tolerance $1.0 \times 10^{-8}$; absolute tolerance $1.0 \times 10^{-10}$; maximum $80$ iterations; relaxation $0.5$ \\
		Linear solver strategy & -- & IPCS blocks: direct MUMPS; SA transport: default solver \\
		\hline
	\end{tabularx}
\end{table}

\subsection{Spatial Verification and Grid Independence}
\label{subsec:app_A_channel_mesh}
To guarantee that the numerical results were independent of the spatial discretization, a formal grid independence study was conducted. Three distinct boundary-layer inflated meshes (Coarse, Medium, and Fine) were evaluated. The geometric properties, node counts, and resolution parameters of this mesh hierarchy are explicitly detailed in Table~\ref{tab:app_A_channel_meshes}. Crucially, the first-layer height was kept constant across all grids to strictly enforce the $y^+ \approx 1$ requirement, while the bulk and inflation-layer resolutions were systematically refined. 

\begin{table}[htbp]
	\centering
	\footnotesize
	\caption{Wall-resolved channel mesh hierarchy for the FEniCS grid-independence study.}
	\label{tab:app_A_channel_meshes}
	\renewcommand{\arraystretch}{1.2}
	\setlength{\tabcolsep}{5pt}
	\begin{tabular*}{\textwidth}{@{\extracolsep{\fill}}lccc@{}}
		\hline
		\textbf{Parameter} & \textbf{Coarse} & \textbf{Medium} & \textbf{Fine} \\
		\hline
		\multicolumn{4}{@{}l}{\textit{Mesh settings}} \\
		Streamwise cells & $10$ & $20$ & $40$ \\
		Wall-normal cells & $42$ & $51$ & $68$ \\
		$\Delta y_1$ [m] & $1.709 \times 10^{-3}$ & $1.709 \times 10^{-3}$ & $1.709 \times 10^{-3}$ \\
		$\Delta y_{\min}$ [m] & $1.709 \times 10^{-3}$ & $1.709 \times 10^{-3}$ & $1.709 \times 10^{-3}$ \\
		$\Delta y_{\max}$ [m] & $7.707 \times 10^{-2}$ & $8.454 \times 10^{-2}$ & $7.589 \times 10^{-2}$ \\
		Inflation layers per wall & $10$ & $14$ & $18$ \\
		Growth ratio & $1.45$ & $1.35$ & $1.25$ \\
		Vertices & $473$ & $1{,}092$ & $2{,}829$ \\
		Elements & $840$ & $2{,}040$ & $5{,}440$ \\
		\hline
		\multicolumn{4}{@{}l}{\textit{Grid metric}} \\
		$U_{\mathrm{bulk}}$ [m/s] & $16.923$ & $18.489$ & $18.472$ \\
		$\epsilon_{\mathrm{rel}}(U_{\mathrm{bulk}})$ [\%] & -- & $9.25$ & $0.09$ \\
		\hline
	\end{tabular*}
	\vspace{2pt}
	\begin{minipage}{0.98\textwidth}
		\scriptsize \textit{Note.} For the $\epsilon_{\mathrm{rel}}$ row, the Medium column gives the Coarse--Medium change and the Fine column gives the Medium--Fine change.
	\end{minipage}
\end{table}

Visual confirmation of this spatial convergence is provided by the wall-normal velocity profiles plotted in Figure~\ref{fig:app_A_channel_mesh_independence}. In each case, the profile is extracted at the streamwise mid-plane, $x=0.5L$, with the wall-normal coordinate $y$ varied across the full channel height. The \texttt{Fine\_FEniCS} grid was consequently selected for the final verification against the Ansys Fluent baseline, while the coarser grids demonstrate that the wall-normal velocity profile is already asymptotically converged.

\begin{figure}[htpb]
	\centering
	\includegraphics[width=\textwidth]{Channel-FEniCS-Mesh-Independence}
	\caption{FEniCS wall-normal velocity profiles at $x=0.5L$ for the Coarse, Medium, and Fine grids.}
	\label{fig:app_A_channel_mesh_independence}
\end{figure}

Furthermore, to ensure a rigorous cross-code validation, an equivalent grid independence study was performed for the commercial Ansys Fluent baseline. The Ansys domain was discretized using a comparable mapped mesh hierarchy, maintaining a strict $y^+ \approx 1$ resolution at the walls. In addition to the residual monitors, the bulk velocity was monitored and required to change by less than $1.0 \times 10^{-3}$ over $250$ iterations before accepting a converged solution. The continuity residual criterion was set to $1.0 \times 10^{-5}$, while the momentum and SA working-variable residual criteria were set to $1.0 \times 10^{-6}$. Mass conservation was also checked from the inlet and outlet mass-flow reports. The properties of these meshes are documented in Table~\ref{tab:app_A_channel_meshes_ansys}. The wall-normal velocity profiles extracted from the Ansys solver at the same mid-channel location, illustrated in Figure~\ref{fig:app_A_channel_mesh_indep_ansys}, exhibited identical asymptotic convergence, confirming that the baseline reference data is mathematically robust and completely independent of spatial discretization errors.

\begin{table}[htbp]
	\centering
	\footnotesize
	\caption{Wall-resolved channel mesh hierarchy used for the Ansys grid-independence study.}
	\label{tab:app_A_channel_meshes_ansys}
	\renewcommand{\arraystretch}{1.2}
	\setlength{\tabcolsep}{5pt}
	\begin{tabular*}{\textwidth}{@{\extracolsep{\fill}}lccc@{}}
		\hline
		\textbf{Parameter} & \textbf{Coarse} & \textbf{Medium} & \textbf{Fine} \\
		\hline
		\multicolumn{4}{@{}l}{\textit{Mesh settings}} \\
		Mesh type & MultiZone & MultiZone & MultiZone \\
		Face type & All quad & All quad & All quad \\
		$N_x$ & $64$ & $160$ & $220$ \\
		$N_y$ & $80$ & $144$ & $180$ \\
		Bias type & Two-sided & Two-sided & Two-sided \\
		Bias factor & $59$ & $28$ & $21$ \\
		Approx. $\Delta y_1$ [m] & $1.700 \times 10^{-3}$ & $1.700 \times 10^{-3}$ & $1.700 \times 10^{-3}$ \\
		Nodes & $7{,}740$ & $26{,}448$ & $43{,}010$ \\
		Elements & $7{,}565$ & $26{,}123$ & $42{,}594$ \\
		\hline
		\multicolumn{4}{@{}l}{\textit{Grid metric}} \\
		$U_{\mathrm{bulk}}$ [m/s] & $18.305$ & $18.114$ & $17.977$ \\
		$\epsilon_{\mathrm{rel}}(U_{\mathrm{bulk}})$ [\%] & -- & $1.04$ & $0.76$ \\
		\hline
	\end{tabular*}
	\vspace{2pt}
	\begin{minipage}{0.98\textwidth}
		\scriptsize \textit{Note.} $N_x$ and $N_y$ are the streamwise and wall-normal edge divisions. For the $\epsilon_{\mathrm{rel}}$ row, the Medium column gives the Coarse--Medium change and the Fine column gives the Medium--Fine change.
	\end{minipage}
\end{table}

\begin{figure}[htpb]
	\centering
	\includegraphics[width=\textwidth]{Channel-Ansys-Mesh-Independence}
	\caption{Ansys wall-normal velocity profiles at $x=0.5L$ for the Coarse, Medium, and Fine grids.}
	\label{fig:app_A_channel_mesh_indep_ansys}
\end{figure}

\FloatBarrier

\section{Verification Benchmark 2: U-Bend Geometry}
\label{sec:app_A_ubend}

\subsection{Physical and Geometric Parameters}
\label{subsec:app_A_ubend_physical}
The second verification benchmark evaluates a highly convective U-bend. Because this FEniCS domain is a 2D planar approximation of the symmetry mid-plane of the 3D Ansys VMFL048 case, the 3D physical pipe diameter ($D$) is explicitly retained as the characteristic length scale. This ensures the Reynolds ($Re_D$) and Dean ($De$) numbers remain non-dimensionally consistent with the commercial baseline. The straight inlet and outlet legs are shorter than in the original VMFL048 Ansys mesh to reduce meshing and computational effort; both the FEniCS and Ansys runs use this same short-leg U-bend geometry. The active FEniCS setup is summarized from \texttt{ConfigUBend\_SpalartAllmaras\_Steady.py}, and the configuration is detailed in Table~\ref{tab:app_A_ubend_physical}.

\begin{table}[htbp]
	\centering
	\footnotesize
	\caption{Geometric, physical, and boundary-condition parameters for the U-bend benchmark.}
	\label{tab:app_A_ubend_physical}
	\renewcommand{\arraystretch}{1.2}
	\begin{tabularx}{\textwidth}{l X}
		\hline
		\textbf{Parameter} & \textbf{U-bend setting} \\
		\hline
		Benchmark type & Short-leg 2D symmetry-plane U-bend based on Ansys VMFL048 \\
		Domain dimensions & Pipe diameter $D=0.028$ m; centreline bend radius $R_c=0.125$ m; straight-leg length $30D=0.840$ m \\
		Characteristic length & Physical pipe diameter $D=0.028$ m \\
		Reference / inlet velocity & Uniform prescribed inlet velocity $\mathbf{u}_{in}=(0,-1.42)$ m/s \\
		Kinematic viscosity & $\nu = 8.9 \times 10^{-7}$ m$^2$/s \\
		Reynolds number & $Re_D \approx 4.47 \times 10^4$ \\
		Dean number & $De \approx 1.50 \times 10^4$ \\
		Body force & $\mathbf{f}=(0,0)$ m/s$^2$ \\
		Driving condition & Prescribed inlet velocity and zero relative static pressure at the outlet \\
		Velocity boundary conditions & No-slip on all curved and straight pipe walls; natural velocity condition at pressure outlet \\
		Pressure boundary conditions & $p=0$ Pa at outlet; pressure left free at inlet and walls \\
		SA boundary conditions & $\tilde{\nu}_{in}=9.34 \times 10^{-5}$ m$^2$/s at inlet; $\tilde{\nu}=0$ at walls; natural condition at outlet \\
			Initial fields & $\mathbf{u}_0=(0,-1.42)$ m/s, $p_0=0$ Pa, $\tilde{\nu}_0=\tilde{\nu}_{in}$ \\
			Verification mesh & \texttt{Medium\_WallResolved}; first layer $\Delta y_1 = 1.20 \times 10^{-5}$ m \\
		\hline
	\end{tabularx}
\end{table}

\subsection{Numerical Simulation Parameters}
\label{subsec:app_A_ubend_numerical}
Due to the adverse pressure gradients and secondary flow separation induced by the streamline curvature, the U-bend represents a significantly more demanding numerical challenge than the channel. Consequently, the FEniCS solver required a smaller pseudo-time step, tighter under-relaxation, and advanced iterative preconditioners to maintain stability, as detailed in Table~\ref{tab:app_A_ubend_numerical}.

\begin{table}[htbp]
	\centering
	\footnotesize
	\caption{Numerical simulation parameters for the highly convective FEniCS U-bend benchmark.}
	\label{tab:app_A_ubend_numerical}
	\renewcommand{\arraystretch}{1.2}
	\begin{tabularx}{\textwidth}{l c X}
		\hline
		\textbf{Parameter} & \textbf{Symbol} & \textbf{Setting} \\
		\hline
		Velocity space & $\mathbf{u}$ & Vector CG2 \\
		Pressure space & $p$ & Scalar CG1 \\
		SA working-variable space & $\tilde{\nu}$ & Scalar CG1 \\
		Quadrature degree & -- & $4$ \\
			Maximum outer Picard steps & $N_{\mathrm{Picard,max}}$ & $1000$ \\
		SA solves per Picard step & -- & $1$ \\
		Outer velocity tolerance & $\epsilon_{\mathbf{u}}$ & $1.0 \times 10^{-4}$ \\
		Outer pressure tolerance & $\epsilon_p$ & $1.0 \times 10^{-4}$ \\
		Outer SA tolerance & $\epsilon_{\tilde{\nu}}$ & $1.0 \times 10^{-6}$ \\
			IPCS pseudo-time step & $\Delta t$ & $1.0 \times 10^{-4}$ s \\
				Maximum IPCS steps per Picard step & -- & $200$ \\
			Maximum IPCS steps for final flow & -- & $600$ \\
			IPCS velocity tolerance & -- & $1.0 \times 10^{-4}$ \\
			IPCS pressure tolerance & -- & $2.0 \times 10^{-3}$ \\
			Pressure-drop convergence monitor & -- & Final-window length $100$ with relative tolerance $5.0\times10^{-4}$; flow convergence required \\
			Velocity relaxation factor & $\alpha_u$ & $0.3$ \\
			Pressure relaxation factor & $\alpha_p$ & $0.1$ \\
			SA relaxation factor & $\alpha_{\tilde{\nu}}$ & $0.01$ \\
			SA pseudo-time stabilization & -- & SA transport pseudo-time step $5.0\times10^{-2}$ s \\
			SA lower bound & $\tilde{\nu}_{\min}$ & $1.0 \times 10^{-12}$ \\
			Wall-distance equation & -- & Relaxed Eikonal, $\sigma_w=0.1$, $G_0=20$ \\
			Wall-distance Newton settings & -- & Relative tolerance $1.0 \times 10^{-8}$; absolute tolerance $1.0 \times 10^{-10}$; maximum $80$ iterations; relaxation $0.5$ \\
			Linear solver strategy & -- & IPCS blocks: GMRES with ILU velocity/correction preconditioning and Hypre-AMG pressure preconditioning; SA transport: default solver \\
			Restart handling & -- & Restart from saved state enabled when checkpoint files are present; otherwise the configured initial fields are used \\
		\hline
	\end{tabularx}
\end{table}

\subsection{Selected FEniCS Mesh and Ansys Grid Independence}
\label{subsec:app_A_ubend_mesh}
	Managing the highly stretched, non-orthogonal inflation layers around the curved arcs of the U-bend required custom mesh generation via Gmsh. The current FEniCS comparison setup points to the wall-resolved \texttt{Medium\_WallResolved} grid documented in Table~\ref{tab:app_A_ubend_meshes}. Similar to the channel benchmark, the wall-adjacent element height is kept at $1.20 \times 10^{-5}$ m to satisfy the low-Reynolds SA wall constraints. A separate coarse--medium--fine FEniCS mesh-convergence study is therefore not claimed for this U-bend verification case.

The reported FEniCS static pressure-drop metric is computed from boundary-measure-weighted mean pressures over the inlet and outlet sections, rather than from local point samples. Let $\Gamma_{\mathrm{in}}$ and $\Gamma_{\mathrm{out}}$ denote the inlet and outlet boundaries, and let $p_k$ be the kinematic pressure solved in FEniCS. The physical static pressure drop reported in pascals is
\begin{equation}
	\Delta p_s
	= \rho\left(\overline{p}_{k,\mathrm{in}}-\overline{p}_{k,\mathrm{out}}\right)
	= \rho\left[
	\frac{1}{|\Gamma_{\mathrm{in}}|}\int_{\Gamma_{\mathrm{in}}} p_k \, \mathrm{d}\Gamma
	-
	\frac{1}{|\Gamma_{\mathrm{out}}|}\int_{\Gamma_{\mathrm{out}}} p_k \, \mathrm{d}\Gamma
	\right],
	\label{eq:app_A_ubend_fenics_delta_p}
\end{equation}
where $|\Gamma|=\int_{\Gamma}1\,\mathrm{d}\Gamma$ is the FEniCS boundary measure and $\rho=997$ kg/m$^3$. In the implementation this corresponds to integrating over the marked inlet and outlet facets with \texttt{ds(2)} and \texttt{ds(3)}, respectively. The density factor is required because the IPCS formulation stores pressure in kinematic form, while the benchmark table reports the static pressure drop in Pa. Although the outlet pressure is prescribed as zero relative pressure, the outlet term is still evaluated as a boundary average so that the diagnostic follows the same inlet-minus-outlet definition on every grid.

	\begin{table}[htbp]
		\centering
		\footnotesize
		\caption{Selected wall-resolved FEniCS U-bend mesh used for comparison with the Ansys Fluent verification baseline.}
		\label{tab:app_A_ubend_meshes}
		\renewcommand{\arraystretch}{1.2}
		\begin{tabularx}{0.82\textwidth}{l X}
			\hline
				\textbf{Parameter} & \textbf{\texttt{Medium\_WallResolved} setting} \\
				\hline
				$h_{\mathrm{bulk}}$ [m] & $1.40 \times 10^{-3}$ \\
				$\Delta y_1$ [m] & $1.20 \times 10^{-5}$ \\
				Inflation layers & $14$ \\
				Growth ratio & $1.30$ \\
				$t_{\mathrm{infl}}$ [m] & $1.53 \times 10^{-3}$ \\
				Vertices & $73{,}978$ \\
				Elements & $144{,}952$ \\
			\hline
		\end{tabularx}
	\end{table}

For the independent Ansys Fluent baseline, a separate mesh-independence study was executed in the same shortened U-bend geometry. Using a parallel unstructured meshing strategy with strict $y^+ \approx 1$ boundary layer inflation (documented in Table~\ref{tab:app_A_ubend_meshes_ansys}), the Ansys velocity profiles at the bend exit across the same pipe-diameter line, shown in Figure~\ref{fig:app_A_ubend_mesh_indep_ansys}, were checked for spatial convergence. For each Ansys solve, convergence was accepted only after the scaled continuity, momentum, and Spalart--Allmaras residuals reached $1.0\times10^{-6}$ and the monitored static pressure drop changed by less than a relative tolerance of $5.0\times10^{-4}$ over the final $100$ iterations. The inlet and outlet mass-flow rates were also compared for each grid to verify the global mass balance before accepting the pressure-drop values. The \texttt{Fine\_Ansys} grid is retained as the commercial verification baseline, while the deliberately coarser \texttt{Coarse\_Ansys} level is selected to bracket the mesh-convergence trend.

\begin{table}[htbp]
	\centering
		\footnotesize
		\caption{Wall-resolved U-bend mesh hierarchy used for the Ansys grid-independence study.}
		\label{tab:app_A_ubend_meshes_ansys}
		\renewcommand{\arraystretch}{1.2}
		\setlength{\tabcolsep}{5pt}
		\begin{tabular*}{\textwidth}{@{\extracolsep{\fill}}lccc@{}}
			\hline
			\textbf{Parameter} & \textbf{Coarse} & \textbf{Medium} & \textbf{Fine} \\
			\hline
			\multicolumn{4}{@{}l}{\textit{Mesh settings}} \\
			Mesh type & Quad-dom. & Quad-dom. & Quad-dom. \\
			$h_{\mathrm{face}}$ [m] & $4.00 \times 10^{-3}$ & $2.00 \times 10^{-3}$ & $1.25 \times 10^{-3}$ \\
			Curvature angle & $22^\circ$ & $12^\circ$ & $8^\circ$ \\
			Inflation option & First layer & First layer & First layer \\
			$\Delta y_1$ [m] & $1.20 \times 10^{-5}$ & $1.20 \times 10^{-5}$ & $1.20 \times 10^{-5}$ \\
			Inflation layers & $10$ & $16$ & $22$ \\
			Growth rate & $1.35$ & $1.22$ & $1.16$ \\
			$t_{\mathrm{infl}}$ [m] & $6.55 \times 10^{-4}$ & $1.59 \times 10^{-3}$ & $2.17 \times 10^{-3}$ \\
			Nodes & $14{,}525$ & $47{,}707$ & $106{,}460$ \\
			Elements & $13{,}984$ & $46{,}640$ & $104{,}763$ \\
			\hline
			\multicolumn{4}{@{}l}{\textit{Grid metric}} \\
			Static pressure drop $\Delta p_s$ [Pa] & $814.109$ & $796.804$ & $792.995$ \\
			$\epsilon_{\mathrm{rel}}(\Delta p_s)$ [\%] & -- & $2.13$ & $0.48$ \\
			\hline
		\end{tabular*}
		\vspace{2pt}
		\begin{minipage}{0.98\textwidth}
			\scriptsize \textit{Note.} Quad-dom. = quad-dominant. For the $\epsilon_{\mathrm{rel}}$ row, the Medium column gives the Coarse--Medium change and the Fine column gives the Medium--Fine change.
		\end{minipage}
	\end{table}

\begin{figure}[htpb]
	\centering
	\includegraphics[width=\textwidth]{UBend-Ansys-Mesh-Independence}
	\caption{Ansys U-bend exit velocity profiles across $D=0.028$ m for the Coarse, Medium, and Fine grids.}
	\label{fig:app_A_ubend_mesh_indep_ansys}
\end{figure}

\FloatBarrier
